\documentclass[11pt]{article}
\usepackage{amsmath,amssymb,amsthm,algorithm,algorithmic}
\usepackage[margin=1in]{geometry}

\newtheorem{theorem}{Theorem}
\newtheorem{lemma}[theorem]{Lemma}
\newtheorem{definition}{Definition}

\title{Problem 4: Largest Palindrome Product}
\author{Project Euler Solutions}
\date{}

\begin{document}
\maketitle

\section{Problem Statement}

Find the largest palindrome made from the product of two 3-digit numbers. Formally, determine
\[
M = \max\bigl\{xy : 100 \le x,y \le 999,\; xy \text{ is a palindrome in base 10}\bigr\}.
\]

\section{Definitions}

\begin{definition}
A non-negative integer $n$ is a \emph{palindrome in base 10} if its decimal representation $d_k d_{k-1} \cdots d_1 d_0$ satisfies $d_i = d_{k-i}$ for all $0 \le i \le k$.
\end{definition}

\begin{definition}
A \emph{6-digit palindrome} is an integer $P$ with $100\,000 \le P \le 999\,999$ that is a palindrome. Its decimal form is $\overline{abccba}$ where $a \in \{1,\dots,9\}$ and $b,c \in \{0,\dots,9\}$.
\end{definition}

\section{Mathematical Development}

\begin{theorem}[Structure of 6-Digit Palindromes]\label{thm:div11}
Every 6-digit palindrome $P = \overline{abccba}$ admits the representation
\[
P = 11(9091a + 910b + 100c).
\]
In particular, $11 \mid P$.
\end{theorem}

\begin{proof}
By positional notation,
\[
P = 10^5 a + 10^4 b + 10^3 c + 10^2 c + 10 b + a = 100001a + 10010b + 1100c.
\]
We verify each coefficient is divisible by 11:
\begin{align*}
100001 &= 11 \times 9091, \\
10010  &= 11 \times 910, \\
1100   &= 11 \times 100.
\end{align*}
Therefore $P = 11(9091a + 910b + 100c)$.
\end{proof}

\begin{lemma}[Factor Constraint]\label{lem:factor}
If $P = xy$ where $P$ is a 6-digit palindrome and $x, y$ are positive integers, then $11 \mid x$ or $11 \mid y$.
\end{lemma}

\begin{proof}
By Theorem~\ref{thm:div11}, $11 \mid P$. Since 11 is prime, Euclid's lemma (if $p$ is prime and $p \mid ab$, then $p \mid a$ or $p \mid b$) gives $11 \mid x$ or $11 \mid y$.
\end{proof}

\begin{theorem}[Product Range]\label{thm:range}
If $100 \le x, y \le 999$, then $10\,000 \le xy \le 998\,001$. Moreover, the largest palindromic product is necessarily a 6-digit palindrome.
\end{theorem}

\begin{proof}
The bounds follow from $100^2 = 10\,000$ and $999^2 = 998\,001$. The largest 5-digit palindrome is $99\,999$. A 6-digit palindromic product exists: $143 \times 707 = 101\,101 = \overline{101101}$. Since $101\,101 > 99\,999$, the maximum palindromic product must be a 6-digit palindrome.
\end{proof}

\begin{theorem}[Optimality]\label{thm:opt}
The largest palindrome that is a product of two 3-digit numbers is
\[
M = 906\,609 = 913 \times 993.
\]
\end{theorem}

\begin{proof}
By Theorem~\ref{thm:range}, any maximal palindromic product of two 3-digit numbers must be a 6-digit palindrome. Therefore Lemma~\ref{lem:factor} applies: if $P = xy$, then at least one of the two factors is divisible by 11.

By commutativity of multiplication, every candidate product can thus be written in the form $xy$ with
\[
110 \le x \le y \le 999,\qquad 11 \mid x.
\]
Hence it suffices to search the finite set
\[
\mathcal{C} = \{(x,y) \in \mathbb{Z}^2 : 110 \le x \le y \le 999,\ 11 \mid x\}.
\]

The algorithm below enumerates $\mathcal{C}$ in descending order of $x$ and, for each fixed $x$, in descending order of $y$. Its pruning rule is sound because for fixed $x$ the products $xy$ strictly decrease as $y$ decreases. Consequently, once $xy \le \text{best}$ no smaller value of $y$ can improve the answer, and the first palindromic product found for a fixed $x$ is the largest one with that value of $x$.

A complete execution over $\mathcal{C}$ returns
\[
\text{best} = 906\,609 = 913 \times 993.
\]
We verify that $906\,609$ is a palindrome and that $913 = 11 \times 83$, so the candidate lies in the restricted search space. Since the search is exhaustive over all relevant pairs, no larger palindromic product exists.
\end{proof}

\section{Algorithm}

\begin{algorithm}[H]
\caption{Largest palindrome product of two 3-digit numbers}
\begin{algorithmic}[1]
\STATE $\text{best} \gets 0$
\FOR{$x = 990$ \textbf{downto} $110$ \textbf{step} $-11$}
    \FOR{$y = 999$ \textbf{downto} $x$ \textbf{step} $-1$}
        \STATE $P \gets x \cdot y$
        \IF{$P \le \text{best}$}
            \STATE \textbf{break} \COMMENT{no further $y$ can improve}
        \ENDIF
        \IF{\textsc{IsPalindrome}($P$)}
            \STATE $\text{best} \gets P$; \textbf{break}
        \ENDIF
    \ENDFOR
\ENDFOR
\RETURN best
\end{algorithmic}
\end{algorithm}

\begin{theorem}[Algorithm correctness]
The algorithm returns the largest palindrome of the form $xy$ with $100 \le x,y \le 999$.
\end{theorem}

\begin{proof}
By Theorem~\ref{thm:range} and Lemma~\ref{lem:factor}, every optimal pair may be reordered so that it belongs to
\[
\mathcal{C} = \{(x,y) : 110 \le x \le y \le 999,\ 11 \mid x\}.
\]
The outer loop enumerates every admissible value of $x$ in $\mathcal{C}$ exactly once, and for each such $x$ the inner loop enumerates admissible $y$ values from $999$ down to $x$.

We justify the early exits.
\begin{enumerate}
\item If $P \le \text{best}$, then for every later inner-loop value $y' < y$ we have $xy' < xy \le \text{best}$, so no skipped pair can improve the answer.
\item If $P$ is palindromic, the algorithm sets $\text{best} \gets P$ and breaks. This is valid because any later value satisfies $y' < y$, hence $xy' < xy = P$, so no later pair with the same $x$ can produce a larger palindrome.
\end{enumerate}

Therefore the algorithm never discards a pair that could beat the final answer, and it eventually records the maximum palindromic product.
\end{proof}

\section{Complexity Analysis}

\begin{theorem}[Time Complexity]
The algorithm runs in $O(N^2/11)$ time in the worst case, where $N = 900$.
\end{theorem}

\begin{proof}
The outer loop visits the 3-digit multiples of 11
\[
110, 121, 132, \ldots, 990,
\]
which form an arithmetic progression with
\[
\frac{990 - 110}{11} + 1 = 81
\]
terms. For each outer-loop value, the inner loop performs at most $900$ iterations in the worst case. Each iteration performs $O(1)$ work: one multiplication, one comparison, and at most one palindrome test on a decimal string of length at most $6$. Hence the total work is bounded by
\[
81 \times 900 = 72\,900 = O(N^2/11),
\]
where $N = 900$ is the number of 3-digit integers. Early termination can only decrease this bound.
\end{proof}

\textbf{Space:} $O(1)$. Only a constant number of variables are used.

\section{Answer}

\[
\boxed{906609}
\]

\end{document}
